% Options for packages loaded elsewhere
\PassOptionsToPackage{unicode}{hyperref}
\PassOptionsToPackage{hyphens}{url}
\documentclass[
  11pt,
]{article}
\usepackage{xcolor}
\usepackage[margin=1in]{geometry}
\usepackage{amsmath,amssymb}
\setcounter{secnumdepth}{5}
\usepackage{iftex}
\ifPDFTeX
  \usepackage[T1]{fontenc}
  \usepackage[utf8]{inputenc}
  \usepackage{textcomp} % provide euro and other symbols
\else % if luatex or xetex
  \usepackage{unicode-math} % this also loads fontspec
  \defaultfontfeatures{Scale=MatchLowercase}
  \defaultfontfeatures[\rmfamily]{Ligatures=TeX,Scale=1}
\fi
\usepackage{lmodern}
\ifPDFTeX\else
  % xetex/luatex font selection
\fi
% Use upquote if available, for straight quotes in verbatim environments
\IfFileExists{upquote.sty}{\usepackage{upquote}}{}
\IfFileExists{microtype.sty}{% use microtype if available
  \usepackage[]{microtype}
  \UseMicrotypeSet[protrusion]{basicmath} % disable protrusion for tt fonts
}{}
\makeatletter
\@ifundefined{KOMAClassName}{% if non-KOMA class
  \IfFileExists{parskip.sty}{%
    \usepackage{parskip}
  }{% else
    \setlength{\parindent}{0pt}
    \setlength{\parskip}{6pt plus 2pt minus 1pt}}
}{% if KOMA class
  \KOMAoptions{parskip=half}}
\makeatother
\usepackage{color}
\usepackage{fancyvrb}
\newcommand{\VerbBar}{|}
\newcommand{\VERB}{\Verb[commandchars=\\\{\}]}
\DefineVerbatimEnvironment{Highlighting}{Verbatim}{commandchars=\\\{\}}
% Add ',fontsize=\small' for more characters per line
\usepackage{framed}
\definecolor{shadecolor}{RGB}{248,248,248}
\newenvironment{Shaded}{\begin{snugshade}}{\end{snugshade}}
\newcommand{\AlertTok}[1]{\textcolor[rgb]{0.94,0.16,0.16}{#1}}
\newcommand{\AnnotationTok}[1]{\textcolor[rgb]{0.56,0.35,0.01}{\textbf{\textit{#1}}}}
\newcommand{\AttributeTok}[1]{\textcolor[rgb]{0.13,0.29,0.53}{#1}}
\newcommand{\BaseNTok}[1]{\textcolor[rgb]{0.00,0.00,0.81}{#1}}
\newcommand{\BuiltInTok}[1]{#1}
\newcommand{\CharTok}[1]{\textcolor[rgb]{0.31,0.60,0.02}{#1}}
\newcommand{\CommentTok}[1]{\textcolor[rgb]{0.56,0.35,0.01}{\textit{#1}}}
\newcommand{\CommentVarTok}[1]{\textcolor[rgb]{0.56,0.35,0.01}{\textbf{\textit{#1}}}}
\newcommand{\ConstantTok}[1]{\textcolor[rgb]{0.56,0.35,0.01}{#1}}
\newcommand{\ControlFlowTok}[1]{\textcolor[rgb]{0.13,0.29,0.53}{\textbf{#1}}}
\newcommand{\DataTypeTok}[1]{\textcolor[rgb]{0.13,0.29,0.53}{#1}}
\newcommand{\DecValTok}[1]{\textcolor[rgb]{0.00,0.00,0.81}{#1}}
\newcommand{\DocumentationTok}[1]{\textcolor[rgb]{0.56,0.35,0.01}{\textbf{\textit{#1}}}}
\newcommand{\ErrorTok}[1]{\textcolor[rgb]{0.64,0.00,0.00}{\textbf{#1}}}
\newcommand{\ExtensionTok}[1]{#1}
\newcommand{\FloatTok}[1]{\textcolor[rgb]{0.00,0.00,0.81}{#1}}
\newcommand{\FunctionTok}[1]{\textcolor[rgb]{0.13,0.29,0.53}{\textbf{#1}}}
\newcommand{\ImportTok}[1]{#1}
\newcommand{\InformationTok}[1]{\textcolor[rgb]{0.56,0.35,0.01}{\textbf{\textit{#1}}}}
\newcommand{\KeywordTok}[1]{\textcolor[rgb]{0.13,0.29,0.53}{\textbf{#1}}}
\newcommand{\NormalTok}[1]{#1}
\newcommand{\OperatorTok}[1]{\textcolor[rgb]{0.81,0.36,0.00}{\textbf{#1}}}
\newcommand{\OtherTok}[1]{\textcolor[rgb]{0.56,0.35,0.01}{#1}}
\newcommand{\PreprocessorTok}[1]{\textcolor[rgb]{0.56,0.35,0.01}{\textit{#1}}}
\newcommand{\RegionMarkerTok}[1]{#1}
\newcommand{\SpecialCharTok}[1]{\textcolor[rgb]{0.81,0.36,0.00}{\textbf{#1}}}
\newcommand{\SpecialStringTok}[1]{\textcolor[rgb]{0.31,0.60,0.02}{#1}}
\newcommand{\StringTok}[1]{\textcolor[rgb]{0.31,0.60,0.02}{#1}}
\newcommand{\VariableTok}[1]{\textcolor[rgb]{0.00,0.00,0.00}{#1}}
\newcommand{\VerbatimStringTok}[1]{\textcolor[rgb]{0.31,0.60,0.02}{#1}}
\newcommand{\WarningTok}[1]{\textcolor[rgb]{0.56,0.35,0.01}{\textbf{\textit{#1}}}}
\usepackage{longtable,booktabs,array}
\newcounter{none} % for unnumbered tables
\usepackage{calc} % for calculating minipage widths
% Correct order of tables after \paragraph or \subparagraph
\usepackage{etoolbox}
\makeatletter
\patchcmd\longtable{\par}{\if@noskipsec\mbox{}\fi\par}{}{}
\makeatother
% Allow footnotes in longtable head/foot
\IfFileExists{footnotehyper.sty}{\usepackage{footnotehyper}}{\usepackage{footnote}}
\makesavenoteenv{longtable}
\usepackage{graphicx}
\makeatletter
\newsavebox\pandoc@box
\newcommand*\pandocbounded[1]{% scales image to fit in text height/width
  \sbox\pandoc@box{#1}%
  \Gscale@div\@tempa{\textheight}{\dimexpr\ht\pandoc@box+\dp\pandoc@box\relax}%
  \Gscale@div\@tempb{\linewidth}{\wd\pandoc@box}%
  \ifdim\@tempb\p@<\@tempa\p@\let\@tempa\@tempb\fi% select the smaller of both
  \ifdim\@tempa\p@<\p@\scalebox{\@tempa}{\usebox\pandoc@box}%
  \else\usebox{\pandoc@box}%
  \fi%
}
% Set default figure placement to htbp
\def\fps@figure{htbp}
\makeatother
\setlength{\emergencystretch}{3em} % prevent overfull lines
\providecommand{\tightlist}{%
  \setlength{\itemsep}{0pt}\setlength{\parskip}{0pt}}
% Shared LaTeX preamble for all Data Mining / ISLR chapter documents.
% Provides \makecoverpage{<assignment title>} — a standalone title page
% with fixed student identity, so every chapter/lab looks uniform.

\usepackage{graphicx}

\newcommand{\makecoverpage}[1]{%
\begin{titlepage}
\centering
\vspace*{2cm}

{\Large \textbf{Adventist University of Central Africa}}\\[0.4cm]
{\large Master's in Big Data Analytics}\\[2cm]

{\Huge \textbf{#1}}\\[0.5cm]
{\large \textit{Introduction to Statistical Learning with Applications in R}}\\[3cm]

\begin{tabular}{rl}
\textbf{Programme:} & MSc Big Data Analytics \\
\textbf{Course Title:} & Data Mining \\
\textbf{Course Code:} & MSDA 9124 \\
\textbf{Lecturer:} & Dr. Pacifique Nizeyimana \\
\textbf{Student Name:} & Wayne Rubangisa \\
\textbf{Student ID:} & 101220 \\
\end{tabular}

\vfill

{\large September 2026}

\end{titlepage}
}

\usepackage{amsmath}
\usepackage{bookmark}
\IfFileExists{xurl.sty}{\usepackage{xurl}}{} % add URL line breaks if available
\urlstyle{same}
\hypersetup{
  hidelinks,
  pdfcreator={LaTeX via pandoc}}

\author{}
\date{\vspace{-2.5em}}

\begin{document}

\makecoverpage{Chapter 3 --- Linear Regression \\ Exercises 1, 6, 11}

\tableofcontents
\newpage

\section{Conceptual Exercises}\label{conceptual-exercises}

\subsection{Exercise 1 --- Null Hypotheses in Table
3.4}\label{exercise-1-null-hypotheses-in-table-3.4}

\textbf{Question.} Describe the null hypotheses to which the p-values
given in Table 3.4 correspond. Explain what conclusions you can draw
based on these p-values. Your explanation should be phrased in terms of
\texttt{sales}, \texttt{TV}, \texttt{radio}, and \texttt{newspaper},
rather than in terms of the coefficients of the linear model.

\textbf{Setting up the model.} The Advertising data are fit with a
single multiple regression model in which \texttt{sales} is regressed
simultaneously on all three advertising budgets:

\[
\text{sales} = \beta_0 + \beta_1 \cdot \text{TV} + \beta_2 \cdot \text{radio} + \beta_3 \cdot \text{newspaper} + \epsilon .
\]

Table 3.4 (reproduced below from the text) reports the estimated
coefficient, its standard error, the resulting \(t\)-statistic, and the
p-value for each term:

{\def\LTcaptype{none} % do not increment counter
\begin{longtable}[]{@{}lrrrl@{}}
\toprule\noalign{}
Predictor & Coefficient & Std. error & \(t\)-statistic & p-value \\
\midrule\noalign{}
\endhead
\bottomrule\noalign{}
\endlastfoot
Intercept & 2.939 & 0.3119 & 9.42 & \textless{} 0.0001 \\
TV & 0.046 & 0.0014 & 32.81 & \textless{} 0.0001 \\
radio & 0.189 & 0.0086 & 21.89 & \textless{} 0.0001 \\
newspaper & \(-0.001\) & 0.0059 & \(-0.18\) & 0.8599 \\
\end{longtable}
}

\textbf{Plain-language explanation.} Each p-value in the table answers
one narrow question: \emph{if I already know how much was spent on the
other two advertising channels, does knowing the spend on this
particular channel still help me predict \texttt{sales}?} This is
different from asking whether the channel is correlated with
\texttt{sales} on its own --- it is asking whether the channel carries
\textbf{extra}, non-redundant information about \texttt{sales} once the
other two channels are held fixed.

\begin{itemize}
\tightlist
\item
  For \textbf{TV}, the p-value is tiny (\texttt{\textless{}\ 0.0001}).
  This says: it is essentially impossible that TV spending has \emph{no}
  effect on \texttt{sales} once radio and newspaper spending are held
  constant --- the data would be extremely surprising under that
  assumption. So we conclude TV advertising is genuinely associated with
  higher sales, over and above whatever radio and newspaper are doing.
\item
  For \textbf{radio}, the same reasoning applies: the p-value is tiny,
  so radio spending appears to have a real, additional effect on sales,
  holding TV and newspaper fixed.
\item
  For \textbf{newspaper}, the p-value is large (\texttt{0.8599}). This
  says the data are entirely consistent with newspaper spending having
  \emph{no} additional effect on sales once TV and radio are already
  accounted for. We do \textbf{not} have evidence that spending more on
  newspaper advertising increases sales, after controlling for the other
  two channels.
\end{itemize}

\textbf{Formal statement.} For each coefficient \(\beta_j\),
\(j \in \{\text{TV},
\text{radio}, \text{newspaper}\}\), the p-value corresponds to testing

\[
H_0 : \beta_j = 0 \qquad \text{vs.} \qquad H_a : \beta_j \neq 0,
\]

holding the other two predictors fixed. At any conventional significance
level (e.g.~\(\alpha = 0.05\)), we \textbf{reject} \(H_0\) for TV and
radio (strong evidence of association with \texttt{sales}), and we
\textbf{fail to reject} \(H_0\) for newspaper (no evidence of
association with \texttt{sales}, once TV and radio are already in the
model).

\textbf{A subtlety worth flagging.} This does \emph{not} mean newspaper
spending is unrelated to sales in a simple, one-predictor regression ---
on its own it is mildly positively correlated with \texttt{sales}. The
reason it drops out here is that newspaper spending is itself correlated
with radio spending (markets that spend more on radio also tend to spend
more on newspaper). Once radio is in the model, it ``absorbs'' the
association that newspaper appeared to have on its own, and newspaper
has nothing left to contribute. This is a classic confounding pattern,
not a contradiction.

\newpage

\subsection{\texorpdfstring{Exercise 6 --- The Least Squares Line Passes
Through
\((\bar x,\bar y)\)}{Exercise 6 --- The Least Squares Line Passes Through (\textbackslash bar x,\textbackslash bar y)}}\label{exercise-6-the-least-squares-line-passes-through-bar-xbar-y}

\textbf{Question.} Using (3.4), argue that in the case of simple linear
regression, the least squares line always passes through the point
\((\bar x, \bar y)\).

\textbf{Plain-language idea.} The least squares line is chosen to
minimize the total squared vertical distance between the data and the
line. It turns out that the intercept term is \emph{defined} in exactly
the way needed to force the line through the ``average point'' of the
data --- the point whose coordinates are the mean of \(x\) and the mean
of \(y\). So this isn't a coincidence that happens to hold for some
datasets; it's built into the formula for \(\hat\beta_0\) itself, and
therefore holds for \emph{every} dataset.

\textbf{Formal argument.} Equation (3.4) gives the least squares
estimates for simple linear regression as

\[
\hat\beta_1 = \frac{\sum_{i=1}^n (x_i - \bar x)(y_i - \bar y)}{\sum_{i=1}^n (x_i - \bar x)^2},
\qquad
\hat\beta_0 = \bar y - \hat\beta_1 \bar x .
\]

The fitted regression line is \(\hat y = \hat\beta_0 + \hat\beta_1 x\).
Evaluate this line at \(x = \bar x\):

\[
\hat y \Big|_{x = \bar x} \;=\; \hat\beta_0 + \hat\beta_1 \bar x
\;=\; \big(\bar y - \hat\beta_1 \bar x\big) + \hat\beta_1 \bar x
\;=\; \bar y .
\]

The \(\hat\beta_1 \bar x\) terms cancel exactly, leaving
\(\hat y = \bar y\) when \(x = \bar x\). Hence the point
\((\bar x, \bar y)\) always satisfies the fitted equation --- that is,
it always lies on the least squares line, regardless of the particular
data values. \(\blacksquare\)

\textbf{Numerical check.} The algebra above holds for \emph{any} numeric
data, so the following chunk checks it on one arbitrary, fixed data set.
Because the data are fixed constants (no random number generation
involved), the result is fully deterministic:

\begin{Shaded}
\begin{Highlighting}[]
\NormalTok{x }\OtherTok{\textless{}{-}} \FunctionTok{c}\NormalTok{(}\DecValTok{1}\NormalTok{, }\DecValTok{2}\NormalTok{, }\DecValTok{3}\NormalTok{, }\DecValTok{4}\NormalTok{, }\DecValTok{5}\NormalTok{, }\DecValTok{6}\NormalTok{, }\DecValTok{7}\NormalTok{)}
\NormalTok{y }\OtherTok{\textless{}{-}} \FunctionTok{c}\NormalTok{(}\FloatTok{2.1}\NormalTok{, }\FloatTok{3.4}\NormalTok{, }\FloatTok{5.0}\NormalTok{, }\FloatTok{4.8}\NormalTok{, }\FloatTok{6.3}\NormalTok{, }\FloatTok{7.1}\NormalTok{, }\FloatTok{7.9}\NormalTok{)}

\NormalTok{fit }\OtherTok{\textless{}{-}} \FunctionTok{lm}\NormalTok{(y }\SpecialCharTok{\textasciitilde{}}\NormalTok{ x)}

\NormalTok{xbar }\OtherTok{\textless{}{-}} \FunctionTok{mean}\NormalTok{(x)}
\NormalTok{ybar }\OtherTok{\textless{}{-}} \FunctionTok{mean}\NormalTok{(y)}
\NormalTok{y\_hat\_at\_xbar }\OtherTok{\textless{}{-}} \FunctionTok{as.numeric}\NormalTok{(}\FunctionTok{predict}\NormalTok{(fit, }\AttributeTok{newdata =} \FunctionTok{data.frame}\NormalTok{(}\AttributeTok{x =}\NormalTok{ xbar)))}

\FunctionTok{data.frame}\NormalTok{(}
  \AttributeTok{quantity =} \FunctionTok{c}\NormalTok{(}\StringTok{"mean(x)"}\NormalTok{, }\StringTok{"mean(y)"}\NormalTok{, }\StringTok{"fitted value at mean(x)"}\NormalTok{),}
  \AttributeTok{value    =} \FunctionTok{c}\NormalTok{(xbar, ybar, y\_hat\_at\_xbar)}
\NormalTok{)}
\end{Highlighting}
\end{Shaded}

\begin{verbatim}
##                  quantity    value
## 1                 mean(x) 4.000000
## 2                 mean(y) 5.228571
## 3 fitted value at mean(x) 5.228571
\end{verbatim}

\begin{Shaded}
\begin{Highlighting}[]
\FunctionTok{isTRUE}\NormalTok{(}\FunctionTok{all.equal}\NormalTok{(y\_hat\_at\_xbar, ybar))}
\end{Highlighting}
\end{Shaded}

\begin{verbatim}
## [1] TRUE
\end{verbatim}

The final line evaluates to \texttt{TRUE}: the fitted value at
\(\bar x\) equals \(\bar y\) exactly (up to floating-point rounding),
confirming the proof above.

\newpage

\section{Applied Exercises}\label{applied-exercises}

\subsection{Exercise 11 --- Regression Through the
Origin}\label{exercise-11-regression-through-the-origin}

\textbf{Question setup.} We generate a predictor \texttt{x} and a
response \texttt{y} as follows, and investigate the \(t\)-statistic for
\(H_0: \beta = 0\) in simple linear regression \textbf{without an
intercept}.

\begin{Shaded}
\begin{Highlighting}[]
\FunctionTok{set.seed}\NormalTok{(}\DecValTok{1}\NormalTok{)}
\NormalTok{x }\OtherTok{\textless{}{-}} \FunctionTok{rnorm}\NormalTok{(}\DecValTok{100}\NormalTok{)}
\NormalTok{y }\OtherTok{\textless{}{-}} \DecValTok{2} \SpecialCharTok{*}\NormalTok{ x }\SpecialCharTok{+} \FunctionTok{rnorm}\NormalTok{(}\DecValTok{100}\NormalTok{)}
\end{Highlighting}
\end{Shaded}

\subsubsection{\texorpdfstring{(a) Regress \texttt{y} onto \texttt{x},
without an
intercept}{(a) Regress y onto x, without an intercept}}\label{a-regress-y-onto-x-without-an-intercept}

\begin{Shaded}
\begin{Highlighting}[]
\NormalTok{fit\_yx }\OtherTok{\textless{}{-}} \FunctionTok{lm}\NormalTok{(y }\SpecialCharTok{\textasciitilde{}}\NormalTok{ x }\SpecialCharTok{+} \DecValTok{0}\NormalTok{)}
\FunctionTok{summary}\NormalTok{(fit\_yx)}
\end{Highlighting}
\end{Shaded}

\begin{verbatim}
## 
## Call:
## lm(formula = y ~ x + 0)
## 
## Residuals:
##     Min      1Q  Median      3Q     Max 
## -1.9154 -0.6472 -0.1771  0.5056  2.3109 
## 
## Coefficients:
##   Estimate Std. Error t value Pr(>|t|)    
## x   1.9939     0.1065   18.73   <2e-16 ***
## ---
## Signif. codes:  0 '***' 0.001 '**' 0.01 '*' 0.05 '.' 0.1 ' ' 1
## 
## Residual standard error: 0.9586 on 99 degrees of freedom
## Multiple R-squared:  0.7798, Adjusted R-squared:  0.7776 
## F-statistic: 350.7 on 1 and 99 DF,  p-value: < 2.2e-16
\end{verbatim}

\textbf{Comment.} Since \texttt{y} was generated as (approximately)
\(2x\) plus noise, we would expect the slope estimate \(\hat\beta\) to
come out close to \(2\), with a standard error that is small relative to
the estimate, and hence a large \(t\)-statistic and a very small p-value
--- strong evidence against \(H_0 : \beta = 0\). Compare these
expectations against the actual \texttt{Estimate}, \texttt{Std.\ Error},
\texttt{t\ value}, and \texttt{Pr(\textgreater{}\textbar{}t\textbar{})}
printed in the \texttt{summary()} output above: the fitted slope should
sit close to \(2\), and the association should come through as highly
statistically significant.

\subsubsection{\texorpdfstring{(b) Regress \texttt{x} onto \texttt{y},
without an
intercept}{(b) Regress x onto y, without an intercept}}\label{b-regress-x-onto-y-without-an-intercept}

\begin{Shaded}
\begin{Highlighting}[]
\NormalTok{fit\_xy }\OtherTok{\textless{}{-}} \FunctionTok{lm}\NormalTok{(x }\SpecialCharTok{\textasciitilde{}}\NormalTok{ y }\SpecialCharTok{+} \DecValTok{0}\NormalTok{)}
\FunctionTok{summary}\NormalTok{(fit\_xy)}
\end{Highlighting}
\end{Shaded}

\begin{verbatim}
## 
## Call:
## lm(formula = x ~ y + 0)
## 
## Residuals:
##     Min      1Q  Median      3Q     Max 
## -0.8699 -0.2368  0.1030  0.2858  0.8938 
## 
## Coefficients:
##   Estimate Std. Error t value Pr(>|t|)    
## y  0.39111    0.02089   18.73   <2e-16 ***
## ---
## Signif. codes:  0 '***' 0.001 '**' 0.01 '*' 0.05 '.' 0.1 ' ' 1
## 
## Residual standard error: 0.4246 on 99 degrees of freedom
## Multiple R-squared:  0.7798, Adjusted R-squared:  0.7776 
## F-statistic: 350.7 on 1 and 99 DF,  p-value: < 2.2e-16
\end{verbatim}

\textbf{Comment.} Now the roles are reversed: we predict \texttt{x} from
\texttt{y}. Since \(y \approx 2x\), inverting the relationship gives
\(x \approx y/2\), so we would expect the slope estimate here to be
roughly \(0.5\) --- noticeably different from the estimate in part (a).
What is more interesting is what does \emph{not} change: as parts
(c)--(e) show, the \(t\)-statistic testing \(H_0 : \beta = 0\) is
identical in both directions, even though the coefficient estimates and
their standard errors differ.

\subsubsection{(c) Relationship between (a) and
(b)}\label{c-relationship-between-a-and-b}

Both regressions are, in effect, testing the same underlying question:
\emph{is there a linear association between \texttt{x} and \texttt{y}?}
The coefficient estimates in (a) and (b) are not reciprocals of one
another and have different standard errors, because they are measured on
different scales (slope of \(y\) on \(x\) versus slope of \(x\) on
\(y\)). However, the strength of evidence for association --- captured
by the \(t\)-statistic and p-value --- turns out to be exactly the same
in both directions. Parts (d) and (e) show why this must be true
algebraically, and the R output above/below confirms it numerically.

\subsubsection{\texorpdfstring{(d) Algebraic derivation of the
\(t\)-statistic
formula}{(d) Algebraic derivation of the t-statistic formula}}\label{d-algebraic-derivation-of-the-t-statistic-formula}

For regression through the origin, (3.38) gives

\[
\hat\beta = \frac{\sum_{i=1}^n x_i y_i}{\sum_{i'=1}^n x_{i'}^2},
\]

and the exercise gives the standard error as

\[
SE(\hat\beta) = \sqrt{\dfrac{\sum_{i=1}^n (y_i - x_i \hat\beta)^2}{(n-1)\sum_{i'=1}^n x_{i'}^2}} .
\]

We want to show that \(t = \hat\beta / SE(\hat\beta)\) can be rewritten
as

\[
t = \frac{\sqrt{n-1}\,\sum_{i=1}^n x_i y_i}{\sqrt{\left(\sum_{i=1}^n x_i^2\right)\left(\sum_{i'=1}^n y_{i'}^2\right) - \left(\sum_{i'=1}^n x_{i'} y_{i'}\right)^2}} .
\]

\textbf{Step 1 --- expand the residual sum of squares.} Write
\(S_{xy} = \sum_i x_i y_i\), \(S_{xx} = \sum_i x_i^2\),
\(S_{yy} = \sum_i y_i^2\). Expanding the square:

\[
\sum_i (y_i - x_i\hat\beta)^2 = \sum_i y_i^2 - 2\hat\beta \sum_i x_i y_i + \hat\beta^2 \sum_i x_i^2
= S_{yy} - 2\hat\beta S_{xy} + \hat\beta^2 S_{xx}.
\]

\textbf{Step 2 --- substitute \(\hat\beta = S_{xy}/S_{xx}\).}

\[
2\hat\beta S_{xy} = \frac{2 S_{xy}^2}{S_{xx}}, \qquad
\hat\beta^2 S_{xx} = \frac{S_{xy}^2}{S_{xx}} .
\]

So

\[
\sum_i (y_i - x_i\hat\beta)^2 = S_{yy} - \frac{2S_{xy}^2}{S_{xx}} + \frac{S_{xy}^2}{S_{xx}}
= S_{yy} - \frac{S_{xy}^2}{S_{xx}}
= \frac{S_{xx}S_{yy} - S_{xy}^2}{S_{xx}} .
\]

\textbf{Step 3 --- assemble the \(t\)-statistic.}

\[
t = \frac{\hat\beta}{SE(\hat\beta)}
= \frac{S_{xy}/S_{xx}}{\sqrt{\dfrac{1}{(n-1)S_{xx}}\left(\dfrac{S_{xx}S_{yy}-S_{xy}^2}{S_{xx}}\right)}}
= \frac{S_{xy}/S_{xx}}{\sqrt{\dfrac{S_{xx}S_{yy}-S_{xy}^2}{(n-1)S_{xx}^2}}} .
\]

Since \(\sqrt{1/S_{xx}^2} = 1/S_{xx}\) (as \(S_{xx} > 0\)), the
\(S_{xx}\) in the denominator of \(\hat\beta\) cancels with the
\(S_{xx}\) pulled out of the square root:

\[
t = \frac{S_{xy}/S_{xx}}{\dfrac{1}{S_{xx}}\sqrt{\dfrac{S_{xx}S_{yy}-S_{xy}^2}{n-1}}}
= \frac{S_{xy}}{\sqrt{\dfrac{S_{xx}S_{yy}-S_{xy}^2}{n-1}}}
= \frac{\sqrt{n-1}\; S_{xy}}{\sqrt{S_{xx}S_{yy}-S_{xy}^2}},
\]

which is exactly the target formula, written back in full as

\[
t = \frac{\sqrt{n-1}\sum_i x_iy_i}{\sqrt{\left(\sum_i x_i^2\right)\left(\sum_i y_i^2\right) - \left(\sum_i x_iy_i\right)^2}}. \qquad \blacksquare
\]

\textbf{Numerical confirmation in R.} The chunk below computes \(t\)
three independent ways --- by hand from \(\hat\beta\) and
\(SE(\hat\beta)\), from the closed-form formula just derived, and by
reading it off \texttt{lm()}'s own summary --- and checks that all three
agree:

\begin{Shaded}
\begin{Highlighting}[]
\NormalTok{n }\OtherTok{\textless{}{-}} \FunctionTok{length}\NormalTok{(x)}

\NormalTok{beta\_hat }\OtherTok{\textless{}{-}} \FunctionTok{sum}\NormalTok{(x }\SpecialCharTok{*}\NormalTok{ y) }\SpecialCharTok{/} \FunctionTok{sum}\NormalTok{(x}\SpecialCharTok{\^{}}\DecValTok{2}\NormalTok{)}
\NormalTok{se\_beta  }\OtherTok{\textless{}{-}} \FunctionTok{sqrt}\NormalTok{(}\FunctionTok{sum}\NormalTok{((y }\SpecialCharTok{{-}}\NormalTok{ x }\SpecialCharTok{*}\NormalTok{ beta\_hat)}\SpecialCharTok{\^{}}\DecValTok{2}\NormalTok{) }\SpecialCharTok{/}\NormalTok{ ((n }\SpecialCharTok{{-}} \DecValTok{1}\NormalTok{) }\SpecialCharTok{*} \FunctionTok{sum}\NormalTok{(x}\SpecialCharTok{\^{}}\DecValTok{2}\NormalTok{)))}
\NormalTok{t\_manual }\OtherTok{\textless{}{-}}\NormalTok{ beta\_hat }\SpecialCharTok{/}\NormalTok{ se\_beta}

\NormalTok{t\_formula }\OtherTok{\textless{}{-}}\NormalTok{ (}\FunctionTok{sqrt}\NormalTok{(n }\SpecialCharTok{{-}} \DecValTok{1}\NormalTok{) }\SpecialCharTok{*} \FunctionTok{sum}\NormalTok{(x }\SpecialCharTok{*}\NormalTok{ y)) }\SpecialCharTok{/}
  \FunctionTok{sqrt}\NormalTok{(}\FunctionTok{sum}\NormalTok{(x}\SpecialCharTok{\^{}}\DecValTok{2}\NormalTok{) }\SpecialCharTok{*} \FunctionTok{sum}\NormalTok{(y}\SpecialCharTok{\^{}}\DecValTok{2}\NormalTok{) }\SpecialCharTok{{-}}\NormalTok{ (}\FunctionTok{sum}\NormalTok{(x }\SpecialCharTok{*}\NormalTok{ y))}\SpecialCharTok{\^{}}\DecValTok{2}\NormalTok{)}

\NormalTok{t\_lm }\OtherTok{\textless{}{-}} \FunctionTok{summary}\NormalTok{(fit\_yx)}\SpecialCharTok{$}\NormalTok{coefficients[}\StringTok{"x"}\NormalTok{, }\StringTok{"t value"}\NormalTok{]}

\FunctionTok{data.frame}\NormalTok{(}
  \AttributeTok{method =} \FunctionTok{c}\NormalTok{(}\StringTok{"by hand (beta\_hat / SE)"}\NormalTok{, }\StringTok{"closed{-}form formula"}\NormalTok{, }\StringTok{"lm() summary"}\NormalTok{),}
  \AttributeTok{t\_stat =} \FunctionTok{c}\NormalTok{(t\_manual, t\_formula, }\FunctionTok{as.numeric}\NormalTok{(t\_lm))}
\NormalTok{)}
\end{Highlighting}
\end{Shaded}

\begin{verbatim}
##                    method   t_stat
## 1 by hand (beta_hat / SE) 18.72593
## 2     closed-form formula 18.72593
## 3            lm() summary 18.72593
\end{verbatim}

\begin{Shaded}
\begin{Highlighting}[]
\FunctionTok{isTRUE}\NormalTok{(}\FunctionTok{all.equal}\NormalTok{(t\_manual, t\_formula)) }\SpecialCharTok{\&\&}
  \FunctionTok{isTRUE}\NormalTok{(}\FunctionTok{all.equal}\NormalTok{(t\_manual, }\FunctionTok{as.numeric}\NormalTok{(t\_lm)))}
\end{Highlighting}
\end{Shaded}

\begin{verbatim}
## [1] TRUE
\end{verbatim}

The final line should evaluate to \texttt{TRUE}, confirming that the
algebra above matches both a from-scratch numerical computation and what
\texttt{lm()} reports.

\subsubsection{\texorpdfstring{(e) The \(t\)-statistic is the same for
\texttt{y} on \texttt{x} as for \texttt{x} on
\texttt{y}}{(e) The t-statistic is the same for y on x as for x on y}}\label{e-the-t-statistic-is-the-same-for-y-on-x-as-for-x-on-y}

Look at the closed-form expression from part (d):

\[
t = \frac{\sqrt{n-1}\sum_i x_i y_i}{\sqrt{\left(\sum_i x_i^2\right)\left(\sum_i y_i^2\right) - \left(\sum_i x_i y_i\right)^2}} .
\]

This formula treats \(x\) and \(y\) almost, but not quite, symmetrically
--- the only way \(x\) and \(y\) enter is through the three summary
quantities \(S_{xy} = \sum_i x_iy_i\), \(S_{xx} = \sum_i x_i^2\), and
\(S_{yy} = \sum_i
y_i^2\). If instead we regress \(x\) onto \(y\) (without an intercept),
the same derivation applies with the roles of \(x\) and \(y\) swapped
throughout, giving

\[
t' = \frac{\sqrt{n-1}\sum_i y_i x_i}{\sqrt{\left(\sum_i y_i^2\right)\left(\sum_i x_i^2\right) - \left(\sum_i y_ix_i\right)^2}} .
\]

But \(\sum_i x_iy_i = \sum_i y_ix_i\) (multiplication of real numbers is
commutative), and the product \(\left(\sum_i y_i^2\right)\left(\sum_i
x_i^2\right)\) is the same regardless of the order in which the two sums
are multiplied. So the numerator and denominator of \(t'\) are literally
the same expressions as in \(t\), just written in a different order:
\(t' = t\). This is exactly the pattern already glimpsed numerically in
parts (a)--(b): the coefficient estimates differ, but the
\(t\)-statistic for \(H_0 : \beta = 0\) does not.

\subsubsection{(f) With an intercept: same check in
R}\label{f-with-an-intercept-same-check-in-r}

\begin{Shaded}
\begin{Highlighting}[]
\NormalTok{fit\_yx\_int }\OtherTok{\textless{}{-}} \FunctionTok{lm}\NormalTok{(y }\SpecialCharTok{\textasciitilde{}}\NormalTok{ x)}
\NormalTok{fit\_xy\_int }\OtherTok{\textless{}{-}} \FunctionTok{lm}\NormalTok{(x }\SpecialCharTok{\textasciitilde{}}\NormalTok{ y)}

\NormalTok{t\_yx\_int }\OtherTok{\textless{}{-}} \FunctionTok{summary}\NormalTok{(fit\_yx\_int)}\SpecialCharTok{$}\NormalTok{coefficients[}\StringTok{"x"}\NormalTok{, }\StringTok{"t value"}\NormalTok{]}
\NormalTok{t\_xy\_int }\OtherTok{\textless{}{-}} \FunctionTok{summary}\NormalTok{(fit\_xy\_int)}\SpecialCharTok{$}\NormalTok{coefficients[}\StringTok{"y"}\NormalTok{, }\StringTok{"t value"}\NormalTok{]}

\FunctionTok{data.frame}\NormalTok{(}
  \AttributeTok{regression =} \FunctionTok{c}\NormalTok{(}\StringTok{"y \textasciitilde{} x"}\NormalTok{, }\StringTok{"x \textasciitilde{} y"}\NormalTok{),}
  \AttributeTok{t\_stat     =} \FunctionTok{c}\NormalTok{(}\FunctionTok{as.numeric}\NormalTok{(t\_yx\_int), }\FunctionTok{as.numeric}\NormalTok{(t\_xy\_int))}
\NormalTok{)}
\end{Highlighting}
\end{Shaded}

\begin{verbatim}
##   regression  t_stat
## 1      y ~ x 18.5556
## 2      x ~ y 18.5556
\end{verbatim}

\begin{Shaded}
\begin{Highlighting}[]
\FunctionTok{isTRUE}\NormalTok{(}\FunctionTok{all.equal}\NormalTok{(}\FunctionTok{as.numeric}\NormalTok{(t\_yx\_int), }\FunctionTok{as.numeric}\NormalTok{(t\_xy\_int)))}
\end{Highlighting}
\end{Shaded}

\begin{verbatim}
## [1] TRUE
\end{verbatim}

The final line should again evaluate to \texttt{TRUE}. This is not a
coincidence special to regression through the origin: for simple linear
regression \textbf{with} an intercept, it is a standard result that the
\(t\)-statistic for testing \(H_0 : \beta_1 = 0\) satisfies

\[
t^2 = \frac{(n-2)\, r_{xy}^2}{1 - r_{xy}^2},
\]

where \(r_{xy}\) is the sample correlation between \(x\) and \(y\).
Because correlation is symmetric (\(r_{xy} = r_{yx}\)), this expression
is identical whether we regress \(y\) on \(x\) or \(x\) on \(y\) --- so
the two regressions must again produce the same \(t\)-statistic, exactly
as the R output above confirms.

\newpage

\end{document}
